\documentclass[instructions.tex]{subfiles}
\begin{document}
 
\chapter{De la contrition et de ses motifs.}
 
\primo La confession sans la contrition n'est pas une confession, mais une illusion ou un sacrilége. La contrition doit renfermer la douleur du passé et le bon propos pour l'avenir. Si je déteste le péché, sans avoir la volonté de le quitter, ou si je quitte le péché sans l'avoir détesté, je n'ai pas la contrition. Pour l'obtenir, demandez-la souvent et humblement à Dieu, et considérez les motifs suivans.

Premier motif. La grandeur et majesté de Dieu. Un sujet qui attaque son roi mérite d'être mis à mort. Que ne méritez-vous pas, vous qui offensez un Dieu devant qui tous les rois sont moins qu'un atome ? Si vous connoissiez la majesté de Dieu, vous mourriez de regret pour l'avoir seulement une fois offensé.

Second motif. La bonté de Dieu. On regarde comme un monstre celui qui, loin d'aimer un bon père, un bon ami, un bienfaiteur, ne leur donne que des marques d'ingratitude. Vous êtes donc pire que tous les monstres, si, loin d'aimer Dieu, vous l'offensez. Il n'y a rien de plus aimable et de plus parfait que Dieu ; il n'y a point de père, d'ami plus tendre et plus bienfaisant : nous ne devrions vivre et respirer que pour l'aimer. Pourquoi ne l'aimez-vous pas. Il ne vous a jamais fait de mal : pourquoi, ingrat, l'outragez-vous ?

Troisiéme motif. La passion de Jésus-Christ. Si, dans un transport de fureur, vous aviez donné la mort à votre père, vous en mourriez de douleur. De quelle affliction ne devriez-vous pas être pénétré, d'avoir donné la mort à Jésus-Christ, votre Sauveur, le meilleur de tous les peres ? Autant de péchés que vous avez commis, ce sont autant de coups que vous lui avez portés, et de plaies que vous lui avez renouvellées. Aimeriez-vous encore le péché, en voyant un Dieu qui pleure, qui souffre et qui meurt pour l'expier ?

Quatriéme motif. L'état de votre âme. Si votre corps était couvert d'une lépre qui vous obligeât à vivre séparé de la société des hommes, ou s'il était défiguré par un ulcère horrible, combien n'en gémiriez-vous pas ! Pourquoi donc n'êtes-vous point touché de voir votre âme séparée de Dieu, défigurée et morte dans le péché ?

Cinquiéme motif. Le ciel perdu. Quand vous perdez un procès, vous en êtes affligé : la perte d'un ami, la mort d'un parent, vous font verser des larmes. Par vos péchés vous avez perdu le ciel, vous avez perdu votre âme, vous avez perdu votre Dieu, et vous n'en êtes point affligé ! O que vous êtes à plaindre dans cette insensibilité !

Sixiéme motif. L'enfer. Vous diriez d'un homme qui serait condamné à être roué ou brûlé vif, qu'il faut qu'il soit bien coupable ; mais si cet homme ne se souciait point de son malheur, vous jugeriez qu'il est fou. Êtes-vous moins coupable, vous qui avez plus de cent fois mérité les feux de l'enfer ? Êtes-vous moins insensé et moins aveugle, vous qui riez, et qui n'êtes point touché d'un état si misérable ?

Entrez en esprit dans ces prisons brûlantes ; écoutez les cris, les lamentations de tant de malheureux qui, au milieu des feux, pleurent inconsolablement leurs péchés et la perte de Dieu, et dites : Voilà ce que j'ai mérité par mes péchés, par mon libertinage et par ma désobéissance à la loi de Dieu, et voilà où je serai bientôt, si je ne change de vie. Oh ! que vous êtes malheureux, si ces considérations ne vous inspirent pas de l'horreur et du regret du péché mortel !

\secundo Quant aux péchés véniels, souvenez-vous, 1. qu'ils affligent Jésus-Christ, qui les a pleurés, et qu'il a plus de douleur d'un seul de vos péchés véniels, que jamais tous les hommes ensemble n'en ont eu de tous les péchés mortels.

2. Que les péchés véniels offensent la sainteté de Dieu, qu'ils refroidissent l'amour que Dieu a pour vous, et refroidissent celui que vous devez avoir pour lui ; qu'ils contristent l'Esprit saint, et que peu à peu ils l'éteignent et l'éloignent de vous.

3. Souvenez-vous enfin que les péchés légers et fréquens vous portent un grand préjudice. Il est vrai qu'ils n'ôtent pas la grâce, ils ne donnent pas la mort à l'ame ; mais ils sont la maladie de l'ame : et de même que les langueurs et les infirmités du corps conduisent peu à peu à la mort, de même les péchés véniels et une vie tiéde conduisent au péché mortel.

Dans vos confessions servez-vous de quelques-unes de ces considérations, pour vous exciter à la douleur de vos fautes.
 
\section{Résolutions}
 
\primo Je demanderai à Dieu la contrition plusieurs jours avant que d'aller à confesse, et je ferai même quelques bonnes œuvres pour l'obtenir plus sûrement.

\secundo Je n'entrerai jamais au tribunal sacré, sans m'être sérieusement excité à la détestation de mes péchés, en me pénétrant vivement des motifs qui me toucheront le plus, parmi ceux que la religion me présente.

\tertio Enfin, si je n'ai que des péchés fort légers à déclarer, afin d'assurer ma contrition, j'en accuserai quelques graves, que je puiserai dans ma vie passée. 

\prayer{Mon Dieu, faites-moi la grâce de m'approcher constamment du sacrement de pénitence avec les dispositions qu'il exige.}
 
\end{document}
